% A29: canonical-class detection on a cyclic prime-power packet.
% Authored in work/; never edits final/.
\subsection{A prime-power canonical-class criterion at a prescribed point}
\label{sec:prime-power-canonical}

\paragraph{Finding A29 (detecting a nonzero primary canonical class).}
Let $C/k$ be a smooth projective geometrically connected genus-two
curve over a finite field $k=\mathbf F_q$ of characteristic $p$, let
$P\in C(k)$, and put
\[
 G=J_C(k),\qquad \kappa_P=[K_C-2P]\in G.
 \tag{PC.1}
\]
Suppose that an odd prime $r\ne p$ satisfies
\[
 \operatorname{ord}_r(p)\ \text{is even},\qquad
 (\kappa_P)_r\ne0,
 \tag{PC.2}
\]
where the subscript denotes the $r$-primary component in the finite
abelian group $G$. Then there are an integer $a\ge1$ and a constant,
geometrically connected cyclic etale $P$-split cover
$D\to C$ of degree $n=r^a$ over $k$ itself such that, for its unique
intermediate quotient $E\to C$ of degree $r^{a-1}$,
\[
 \begin{aligned}
 \mathcal P&=\ker(\operatorname{Nm}:J_D\longrightarrow J_E)^0,\\
 \#\mathcal P(k)&=\frac{\#J_D(k)}{\#J_E(k)},
 \qquad p\nmid\#\mathcal P(k).
 \end{aligned}
 \tag{PC.3}
\]
Consequently the covering formulation of A1 separates the point
section of $P$ from the point section of every distinct $Q\in C(k)$.
The statement only needs the rational $p$-primary group of this
primitive factor to vanish; its geometric $p$-rank need not be zero.
It does not assert that $\#J_D(k)/\#J_C(k)$ is prime to $p$ when
$a>1$.

\paragraph{Recovered inputs and the refinement.}
The prime-degree canonical-class criterion is already proved in
recovered successor \emph{final row A.26}, in every characteristic:
if $-1\in\langle p\rangle\subset(\mathbf Z/r)^\times$ and
$\kappa_P\notin rG$, a prime-degree rational witness exists
\cite[A.26]{ASrun}. Its proof uses the middle-coefficient
Riemann--Roch identity used below. Recovered A1 B.13--B.14 already
identify the primitive factor in a cyclic prime-power tower as the
appropriate rational norm obstruction \cite[B.13--B.14]{A1run}.
The refinement here combines these inputs to replace
$\kappa_P\notin rG$ by the strictly weaker finite-group condition
$(\kappa_P)_r\ne0$. We prove the whole combination, including
its pointed-cover consequence. This attribution does not make a
claim of bibliographic priority beyond the recovered source audit.

\begin{proof}
\emph{1. A cyclic quotient detecting the primary component.}
Decompose the $r$-primary subgroup as
$G_r\simeq\bigoplus_i\mathbf Z/r^{e_i}$.
Since $(\kappa_P)_r\ne0$, at least one of its coordinates is
nonzero. Projection to that coordinate, preceded by the canonical
primary projection of $G$, gives a surjection
\[
 b:G\twoheadrightarrow A=\mathbf Z/r^a,
 \qquad b(\kappa_P)\ne0,
 \tag{PC.4}
\]
with $a=e_i$. Conversely, any cyclic $r$-power quotient detecting
$\kappa_P$ implies $(\kappa_P)_r\ne0$.
In particular (PC.4) does not require $\kappa_P$ to survive in
$G/rG$: an element can be nonzero in a cyclic $r^a$-factor while
being divisible by $r$ in that factor.

Let $d_a=\operatorname{ord}_{r^a}(p)$. Reduction modulo $r$
shows that $\operatorname{ord}_r(p)$ divides $d_a$, so $d_a$ is
even. The group $(\mathbf Z/r^a)^\times$ is cyclic and its
unique element of order two is $-1$. Hence
\[
                  p^{d_a/2}\equiv-1\pmod{r^a}.
 \tag{PC.5}
\]
This uses the actual order at the chosen power. No assertion of
maximal order growth, or non-Wieferich hypothesis, is needed.

\emph{2. The constant pointed Lang cover.}
Write $\Phi$ for $q$-power Frobenius on $J=J_C$ and
$H=\ker b$. The map $\Phi-1$ is an etale isogeny with constant
kernel $J(k)=G$. It factors as
\[
 J\longrightarrow J/H\xrightarrow{\psi}J,
 \qquad \ker\psi=G/H=A.
\]
Pull $\psi$ back by $i_P:C\to J$, $R\mapsto[R-P]$, to obtain
$D=C\times_{J}(J/H)$. This is an etale cover of degree $n$ with
constant deck group $A$. Its fibre over $P$ is the constant kernel
of $\psi$, so it has a chosen rational point $P'$ above $P$ and
the whole fibre is rational.

For completeness, this pullback is geometrically connected.
Otherwise its geometric monodromy would be a proper subgroup
$A_0\subset A$. Quotienting by $A_0$ gives a nontrivial intermediate
etale isogeny $\psi':J'\to J$ whose pullback to $C$ has trivial
geometric monodromy. Over an algebraic closure, $i_P$ therefore
lifts to $v:C\to J'$. Subtract $v(P)$, which lies in $\ker\psi'$.
The Jacobian universal property gives a homomorphism $u:J\to J'$
with $u i_P=v-v(P)$ and $\psi'u=\operatorname{id}_J$.
The homomorphism $u$ is an isogeny, since its composition with
$\psi'$ is an isomorphism and both varieties have dimension two.
Multiplicativity of degree now contradicts $\deg\psi'>1$.
This is the connected Lang-pullback argument of A12.

\emph{3. The genus-two character polynomial and its symmetry.}
Let $\chi:A\to\mu_n$ be a nontrivial characteristic-zero
character. The unramified character Euler product for this cover is
\[
 L_P(\chi,T)
 =\sum_{d\ge0}\left(
   \sum_{Z\in\operatorname{Div}^{+}_d(C)(k)}
      \chi\bigl(b([Z-dP])\bigr)\right)T^d.
\]
The convention for Frobenius translation is the Lang convention
of recovered A1 B.2; reversing that convention replaces every
character by its inverse and leaves the argument unchanged.
In genus two its exact polynomial is
\[
 \begin{aligned}
 L_P(\chi,T)&=1+c_1(\chi)T
                  +q\chi\bigl(b(\kappa_P)\bigr)T^2,\\
 c_1(\chi)&=\sum_{R\in C(k)}\chi\bigl(b([R-P])\bigr).
 \end{aligned}
 \tag{PC.6}
\]
Here is a direct verification of the truncation. For $d>2$,
Riemann--Roch makes the number of effective divisors in each
degree-$d$ line-bundle class independent of the class. Summing
the nontrivial character over $\operatorname{Pic}^d(C)(k)$ gives
zero. In degree two, every line-bundle class has one effective
divisor except the canonical class, which has $q+1$. The excess
is $q\chi(b(\kappa_P))$. The degree-zero coefficient is one,
and degree one gives the displayed point sum. The rational point
$P$ identifies every degree component with $J(k)$ in these sums.

Riemann--Roch in degree one also gives
\[
 c_1(\chi)=\chi\bigl(b(\kappa_P)\bigr)c_1(\chi^{-1}).
 \tag{PC.7}
\]
Indeed $R\mapsto\iota R$, where $R+\iota R\sim K_C$, is a
bijection on $C(k)$. This follows directly because $K_C-R$ has
a unique effective divisor of degree one, defined over $k$.
The identity
$[\iota R-P]=\kappa_P-[R-P]$ proves (PC.7).
Thus no rational Weierstrass point is required.

\emph{4. Every faithful character is a unit at $p$.}
Reduce the cyclotomic integers modulo a fixed prime above $p$,
and denote the resulting character by $\bar\chi$. Since
$p\nmid n$, reduction preserves its order and its faithfulness.
Put $S_\chi=\overline{c_1(\chi)}$. If $j=d_a/2$, (PC.5)
and the integer coefficients of the point sum give
\[
 S_{\chi^{-1}}=S_\chi^{p^j},\qquad
 S_\chi=\bar\chi\bigl(b(\kappa_P)\bigr)S_\chi^{p^j}.
 \tag{PC.8}
\]
Suppose that $L_P(\chi,1)$ vanishes modulo that prime.
Equation (PC.6) gives $S_\chi=-1$, since $q=0$ in the residue
field. Its $p^j$th power is also $-1$, so (PC.8) forces
$\bar\chi(b(\kappa_P))=1$.
For a faithful $\chi$ this contradicts (PC.4). Therefore
\[
 L_P(\chi,1)\ \text{is a unit at }p
       \quad\text{for every character of exact order }r^a.
 \tag{PC.9}
\]
Nonfaithful characters may have zero reduction. The argument
deliberately makes no assertion about those factors.

\emph{5. The top-step Prym has prime-to-$p$ rational order.}
Choose a generator $\sigma$ of $A$ and set
$\tau=\sigma^{n/r}$, of order $r$; let $E=D/\langle\tau\rangle$.
Pullback and norm for $D\to E$ have composition $[r]$ on $J_E$.
They yield a $k$-isogeny
$J_E\times\mathcal P\to J_D$, so finite-field point counts give
the first equality in (PC.3).

For the exact character factorization, the Euler products of
$D/C$ and $E/C$ are products of the corresponding character
Euler factors. The characters for $E/C$ are exactly those of $A$
trivial on $\langle\tau\rangle$. Since $n$ is a prime power,
the remaining characters are exactly the faithful ones. Hence
\[
 \frac{\#J_D(k)}{\#J_E(k)}
       =\prod_{\substack{\chi\in\widehat A\\
                          \operatorname{ord}(\chi)=r^a}}
                       L_P(\chi,1).
 \tag{PC.10}
\]
This also follows by cancelling the imprimitive character factors
in the zeta-function ratios of recovered A1 B.2. The right side
is an integer, by its identification with $\#\mathcal P(k)$.
Each factor is a unit at the chosen prime above $p$ by (PC.9),
so the integer is prime to $p$. Thus
$\mathcal P(k)[p^\infty]=0$. The dimension check is
$\dim\mathcal P=(n-n/r)(g(C)-1)=\varphi(n)$.

\emph{6. The prime-power deck argument separates the point.}
Suppose that a distinct $Q\in C(k)$ collided with $P$ in the
covering formulation of A1. Applied to the pointed cover just
constructed, that formulation supplies $Q'\in D(k)$ above $Q$
with
\[
                  \delta=[Q'-P']\in J_D(k)[p^\infty].
 \tag{PC.11}
\]
This is the same collision-to-cover implication used in recovered
A1 B.1 and in A11; no converse to that implication is used.
The point $(\tau-1)\delta$ is rational, is $p$-primary, and lies in
\[
             \operatorname{im}(\tau-1)=\mathcal P.
 \tag{PC.12}
\]
To verify the equality, the image is connected and lies in the
norm kernel. Conversely, on $\mathcal P$ the endomorphism
$[r]-\sum_{i=0}^{r-1}\tau^i$ is $[r]$ and is divisible by
$\tau-1$ as a polynomial in $\tau$. Multiplication by $r$ is
surjective on $\mathcal P$, proving the reverse inclusion.
Equation (PC.3) therefore gives $(\tau-1)\delta=0$.

It follows that
$\tau Q'+P'-Q'-\tau P'$ is principal. Its four points are
distinct: the deck action is free, and the fibres over $P$ and
$Q$ are disjoint. Its positive and negative parts each have
degree two with simple support. The resulting rational function
gives a separable degree-two map $D_{\bar k}\to\mathbf P^1$;
the simple poles rule out a $p$th power if $p=2$, and the degree
itself excludes inseparability if $p\ne2$.

Since $g(D)=n+1\ge4$, its hyperelliptic pencil is unique.
The order-$r$ automorphism $\tau$ descends to an automorphism
of that projective line of order $r$: the kernel of the descent
action has order at most two. Over $\bar k$ this odd-order
automorphism has a fixed point. Its fibre contains at most two
geometric points, so an order-$r$ cyclic group acts trivially on
that fibre. This gives a fixed point of $\tau$ on $D$, contrary
to its free deck action. The assumed collision is impossible.
\end{proof}

\paragraph{The characteristic-two Weierstrass form and its exact residue.}
If $p=2$ and $W\in C(k)$ is a Weierstrass point, then
$K_C\sim2W$, so $\kappa_P=2[W-P]$. Multiplication by two is
an automorphism on every odd-primary subgroup. Consequently A29
applies precisely when $[W-P]$ has a nonzero $r$-primary component
for some odd prime $r$ with even $\operatorname{ord}_r(2)$.
Any two rational Weierstrass points differ by a class killed by
two, so this condition is independent of their choice.
It supplies a finite-field rational primitive-factor witness with
no size inequality and no analytic hypothesis.

This detector leaves unchanged those points for which the odd part
of $[W-P]$ is supported entirely on primes with odd
$\operatorname{ord}_r(2)$, including the case where $[W-P]$ is
purely $2$-primary. In particular it says nothing at a Weierstrass
point $P$, where $\kappa_P=0$. A constant extension cannot remove
this obstruction to the detector: the map on rational points of
the Jacobian is injective, so the finite order and primary support
of the fixed class $\kappa_P$ do not change. This is a limitation
of the sufficient construction, not a nonexistence theorem for
other cover witnesses. The A09 example already has non-Weierstrass
rational points whose differences from a rational Weierstrass
point are purely $2$-primary, because its rational Jacobian has
order eight.

\paragraph{Strictness on an actual curve.}
The difference from the prime-level canonical detector occurs on
the explicit curve of A09, without an abstract group realization
assumption. Write its point at infinity as $W$, and retain the
point $R=(u,v)\in C(\mathbf F_{16})$ whose class
$L=[R-W]$ has exact order $13$. For the prescribed point $R$,
\[
                   \kappa_R=-2L
 \tag{PC.13}
\]
also has exact order $13$. The finite etale surjective isogeny
$[13]:J_C\to J_C$ has a nonempty fibre over $\kappa_R$.
A point of this finite fibre is rational over some finite
extension $k'/\mathbf F_{16}$; write it as $U$, so that
$13U=\kappa_R$. Thus
$\kappa_R\in13J_C(k')$ while its $13$-primary component remains
nonzero. For any odd prime $r\ne13$, choose an integer $s$ with
$rs\equiv1\pmod{13}$. Then
$\kappa_R=r(s\kappa_R)$ in $J_C(k')$. Consequently
\[
                  \kappa_R\in rJ_C(k')
                    \quad\text{for every odd prime }r.
 \tag{PC.14}
\]
Every prime-level canonical detector in A28 therefore fails over
this field. Nevertheless $\operatorname{ord}_{13}(2)=12$, so
A29 constructs a cyclic $13^a$ pointed witness with the top-step
rational Prym of odd order. Any quotient detecting $\kappa_R$
must have $a\ge2$, because its image in every quotient of order
$13$ is zero. This proves strictness of the canonical detection
criteria on an actual curve. It does not assert that all actual
prime-degree covers fail: the independently constructed geometric
order-$13$ witness of A09 remains available. No explicit degree
for $k'$ is needed or claimed.

Recovered successor A.44 concerns a different pointed-involution
constraint at a Weierstrass base point and its primitive Frobenius
multipliers \cite[A.44]{ASrun}. It is compatible with the present
residual case. In general composite degree an even multiplicative
order alone does not force $-1\in\langle p\rangle$; the use of a
single odd prime power in (PC.5) is essential to the stated
hypothesis and to the simple order-$r$ deck endgame.
